\section{Conclusion}\label{sec:conclusion}

Frozen V-JEPA conditioning benefits a matched low-light restorer in domain and
under the specified leave-one-domain-out protocol. The stronger question is
what the trained branch uses. It tolerates removal of most centred token
variation when its clip mean and a few clip-specific principal directions are
preserved; the tested fixed replacement is much less successful. An exploratory
validation control separates contributions of the clip mean and, conditional
on that mean, adaptive directions; its rank-four versus full-field comparison
remains inconclusive. These results
characterise the trained readout and do not yet provide a smaller encoder or a
validated computational compression recipe. The corrected fixed-rank
training control has an unresolved effect rather than the previously reported
penalty, while the single-seed random control resolves a perceptual benefit
without resolving a fidelity benefit.
Hybrid keys modestly improve LPIPS, spatial grounding does not resolve an improvement, exposure decoupling
degrades the result, and directly matching latent temporal changes worsens
perceptual quality and the measured tLP gap. The study supports measuring
downstream use separately from what a pretrained model can represent, and
making those measurements on converged, reproducible comparisons.
